\section{Related Work}
\label{appendix:relatedwork}
\paragraph{Self-evolving agents.}
Self-evolving agents refer to intelligent agent systems that can continuously absorb and integrate new knowledge, strategies, and tools, and autonomously optimize their capabilities~\citep{fang2025comprehensivesurveyselfevolvingai,gao2025evovlesurvey}. There are currently four main directions of self-evolution~\citep{shao2025misevovle}. The first focuses on model evolution, where agents further refine their parameters through self-generated data, combined with methods such as reinforcement learning~\citep{huang2025rzero,zhou2025safeagent,sun2025seagent}. The second focuses on memory evolution, where agents continuously store experience during interactions and use past experience to guide future interaction~\citep{zhou2025selfchallenging,zhang2025agentexperience,bo2025agentic}. The third focuses on tool evolution, where agents can create, reuse, and optimize tools during interactions, expanding their capability boundaries and improving their proficiency in using tools~\citep{shi2025tool,qu2024exploration}. The fourth focuses on workflow evolution, where agents can also improve performance by optimizing execution pipelines or collaborative processes~\citep{zhang2024aflow,wang2025evoagentx}.
\paragraph{Agent safety.}
Given that the risks posed by LLM agents are not limited to natural language output, the methods for protecting against them differ fundamentally from those for LLM guardrail~\citep{gao2025evovlesurvey,yu2025survey}. Most existing LLM agent protection methods lack generalization ability, TrustAgent~\citep{hua2024trustagent} requires predefined rules, GuardAgent~\citep{xiang2024guardagent} and Conseca~\citep{tsai2025contextual} rely on manually specified trusted contexts. Although some works have begun to focus on adaptation~\citep{zhou2025safeagent,luo2025agrail}, it is still insufficient for LLM agents with diverse output modes and security requirements. Our DEFEND framework, which empowers agents to evolve on their own, will be key to bridging the problem.